\documentclass[letterpaper,fleqn]{article}
\usepackage[spanish,es-noshorthands]{babel}
\usepackage[utf8]{inputenc} 
\usepackage[left=1cm, right=1cm, top=1.5cm, bottom=1.7cm]{geometry}
\usepackage{mathexam}
\usepackage{amsmath}
\usepackage{graphicx}

\ExamClass{\includegraphics[height=16pt]{Images/logo-sed.png} Aritmética $6^{\circ}$}
\ExamName{Fracciones}
\ExamHead{\includegraphics[height=16pt]{Images/logo-colegio.png} IEDAB}
\newcommand{\LineaNombre}{%
\par
\vspace{\baselineskip}
Nombre:\hrulefill \; Curso: \underline{\hspace*{48pt}} \; Fecha: \underline{\hspace*{2.5cm}} \relax
\par}
\let\ds\displaystyle

\begin{document}
\ExamInstrBox{
Respuesta sin justificar mediante procedimiento no será tenida en cuenta en la calificación. Escriba sus respuestas en el espacio indicado. Tiene 45 minutos para contestar esta prueba.}
\LineaNombre
\begin{enumerate}
 \item La señora de Rojas hizo una torta que dividió en 8 partes iguales, de las cuales ella comió una porción, su esposo dos porciones, su hijo mayor dos porciones y su hijo menor una porción. La fracción de la torta que comieron la señor de Rojas y su familia es: \noanswer
 \item Un atleta diariamente da 24 vueltas a una pista. Hoy, cuando corría, sufrió una lesión y solamente había hecho 18 vueltas. ¿Qué  fracción de lo que normalmente corre alcanzó a hacer? \noanswer
\item Juan y Pedro deben llevar cemento para hacer una obra. si Juan lleva $\frac{3}{4}$ de bulto y Pedro $\frac{3}{5}$ de bulto, ¿llevan ambos la misma cantidad?\noanswer
\item Valentina y Reinel comen torta. Si Valentina come $\frac{4}{8}$ de torta y Reinel $\frac{2}{4}$ de torta, ¿comen ambos la misma cantidad de torta?\noanswer
 \end{enumerate}

\end{document}
